\chapter{The Gaussian-query proposal on a growing mesh}\label{u:sec:proposal}

In this section, we construct the proposal used in
\cref{u:sec:fine,u:sec:coarse,u:sec:connections}. We first prove the
Gaussian assembly and reconstruction identities in
\cref{u:lem:assembly}. We then identify the law of the complete histories of the rows and prove the quadrature estimate in \cref{u:lem:quad}.
Throughout, we use the scalar inputs of \cref{u:sec:scalar}.

We follow the Gaussian-conditioning construction of
\cite[Section~4]{oai-typical-upper}. Successive Gaussian revelations along
orthogonalized directions occur in the iterative construction of
Bolthausen~\cite{bolthausen-2014}. Montanari used related increments guided by the Parisi process in an
incremental algorithm for approximate message
passing~\cite{montanari-2021}. Here the durations of the cells grow with
time, and the source and replay constructions retain all the auxiliary
coordinates needed for the row identities.

\section{Growing schedule and source law}\label{u:sec:proposal-schedule}

For $T>0$, let $\mathcal C_T=C([0,T],\R)$. We use the path-space
conventions of \cref{u:sec:setup}. We write $\mathrm{Br}_h$ for the law of a Brownian bridge of duration $h>0$,
which is the image of Wiener measure under
$b(u)\mapsto b(u)-(u/h)b(h)$.
Fix $n$, $h>0$, and an integer $N<n$. We define
\begin{equation}\label{u:eq:growing-schedule}
 t_0=0,\qquad h_j=h e^{t_j/4},\qquad
 t_{j+1}=t_j+h_j,\qquad V=t_{N+1}.
\end{equation}
Cell zero is auxiliary. Query $j\in\{1,\ldots,N\}$ uses
\begin{equation}\label{u:eq:query-targets}
 \begin{aligned}
 s_j^2&=\int_{t_{j-1}}^{t_j}r(v)\dd v,&
 \mathcal T_{j,i}&=\frac{m_i(t_j)-m_i(t_{j-1})}{s_j},\\
 b_j&=\frac1{\sqrt{h_j}}\int_{t_j}^{t_{j+1}}\rho(v)p(v)\dd v.
 \end{aligned}
\end{equation}
Here $m_i(t)=m(t,Y_i(t))$. The scalar reference $b_j$ should not be
confused with the drift $b(v,y)$ of the scalar equation. We use the normalized
inner product and norm
\[
 \ipn{u}{v}=n^{-1}u^{\mathsf T}v,\qquad
 \normn{u}^2=\ipn{u}{u}.
\]

We define the unlinked source law $\mu^{\rm src}_0$ as the independent product
of the laws of the following coordinates:
\begin{itemize}
\item $n$ Wiener paths $D_i$ through $V$ and auxiliary Wiener paths
 $B_i^0$ of duration $h_0$;
\item bridges $\mathrm{br}_{j,i}\sim\mathrm{Br}_{h_j}$ for
 $1\le j\le N$ and $1\le i\le n$;
\item independent $e_j,v_j\sim N(0,I_n)$ and $d_j\sim N(0,2/n)$
 for $1\le j\le N$;
\item a full primitive $g^0\sim\gamma_n$, with ambient matrix
 $W_0=W(g^0)$, and a spin $x\sim\Unif_n$.
\end{itemize}
Unused bridges may be retained as independent ancillary coordinates;
the estimates below do not depend on them. We write $\Omega^{\rm src}$
for this product space. When we use primitive disorder, we adjoin the
antisymmetric variable specified below, independently of the other inputs.

On the auxiliary cell, $B_i=B_i^0$. For $j\ge1$, we append cell $j$ by
\begin{equation}\label{u:eq:append-cell}
 B_i(t_j+u)=B_i(t_j)+\frac{u}{\sqrt{h_j}}e_{j,i}
                    +\mathrm{br}_{j,i}(u),\qquad 0\le u\le h_j.
\end{equation}
Almost surely, and simultaneously for all the cells and rows used here,
a bridge is unchanged when we subtract the line through its two endpoints. We define $M_i=\int\rho\dd B_i+\int\sigma\dd D_i$ and solve the
scalar equation of \cref{u:sec:scalar} in each coordinate. Since the weights in these integrals are deterministic and of bounded
variation, the integrals have the causal path interpretation given
there. In particular, we can evaluate the targets in
\eqref{u:eq:query-targets} before we append \eqref{u:eq:append-cell}.

Applying the Gram--Schmidt procedure to the $\mathcal T_j$ in chronological order,
we obtain directions $u_j$ such that $\normn{u_j}=1$; when a residual vanishes,
the choice of direction is deterministic and measurable. More precisely, let $E$ have
rows that are Euclidean orthonormal and span the remaining subspace. If
$E\mathcal T_j\ne0$, we normalize $E\mathcal T_j$; otherwise, we
normalize the first nonzero column of $E$ in the fixed coordinate
order. Let $q$ be the resulting vector, which has Euclidean norm one; the query
direction is then $u_j=\sqrt n E^{\mathsf T}q$. We complete $q$ to an orthonormal basis
by a fixed measurable rule, rotate $E$ by this basis, and retain all but
its first row as the next remaining frame. We do this also when a
residual vanishes.
We define
\[
 U_j=\frac1n\sum_{l\le j}u_lu_l^{\mathsf T},\qquad
 P_j=I-U_j,\qquad f_j=P_je_j.
\]
The assembled matrix and the completion vectors are defined by
\begin{equation}\label{u:eq:assembled-matrix}
 \begin{aligned}
 W&=\frac1n\sum_{j=1}^N
       (u_jf_j^{\mathsf T}+f_ju_j^{\mathsf T})
       +\frac1n\sum_{j=1}^N d_ju_ju_j^{\mathsf T}+P_NW_0P_N,\\
 C_j&=U_je_j+P_jv_j.
 \end{aligned}
\end{equation}
Note that the scalar recursion uses none of $d_j$, $v_j$, $W_0$, or $x$.

\begin{lemma}[Gaussian assembly, completion, and path reconstruction]
\label{u:lem:assembly}
Under $\mu^{\rm src}_0$, conditional on all initial data
$(D,B^0,\mathrm{br})$, the joint law of $(W,(C_j)_j)$ is
$\mathcal G_n\otimes N(0,I_n)^{\otimes N}$. This law does not depend
on the initial data. The fixed-matrix replay endpoint
\begin{equation}\label{u:eq:replay-endpoint}
 e_j^{\rm rep}=P_jWu_j+U_jC_j
\end{equation}
equals $e_j$ on the source realization. The reconstructed $(B_i,D_i)$
have the product of $n$ independent Wiener-pair laws through $V$.
\end{lemma}

\begin{proof}
We first prove the conditional Gaussian law by induction on the
remaining queries. We then expand the assembled matrix to identify the
replay endpoints. Finally, we compute the covariance of each appended
path to show that the reconstructed pairs have the Wiener law.

\paragraph{Conditional Gaussian law.}
For a fixed remaining frame, the column has Gaussian coordinates,
consisting of a diagonal $N(0,2/n)$ entry and independent off-diagonal $N(0,1/n)$
entries. The projection $P_je_j$ gives the off-diagonal column:
when we multiply a Euclidean unit direction by $\sqrt n$ to obtain $u_j$,
this cancels the $1/\sqrt n$ scale of the column. Conditional on the previously
read $e$'s and the initial data, the next frame is fixed.

We decompose the fresh $e_j,v_j$ along $U_j,P_j$. Since the projections are orthogonal, the four Gaussian components
are independent. It follows that $U_je_j+P_jv_j$ is standard Gaussian and independent of $P_je_j$
and $d_j$. By backward induction over the remaining queries, we find that the
unrevealed block is a GOE compression on the remaining subspace, with
the same $1/n$ normalization of the variance. It is independent of the
part already read and of the later completion outputs. The base case is given by the independent
ambient matrix $W_0$. Using orthogonal invariance as in \cref{u:sec:goe}, we assemble
each column with this block. This proves the assertion on the joint law, including
all later $C_j$, and that this law does not depend on the initial data.

\paragraph{Replay endpoints.}
Expanding \eqref{u:eq:assembled-matrix}, we find
$P_jWu_j=P_je_j$. Indeed, the earlier terms lie in the range of $U_j$, and the later terms
and the final compression annihilate $u_j$. We also have
$U_jC_j=U_je_j$, so \eqref{u:eq:replay-endpoint} equals $e_j$.
By induction, we conclude that every replay target and frame agrees
pointwise with the source, as does every initial segment of the
scalar paths.

\paragraph{Pair paths.}
On one cell, $(u/\sqrt h)e+\mathrm{br}(u)$ is a centered continuous
Gaussian path, and its covariance at $u,v$ is
\[
 \frac{uv}{h}+\min(u,v)-\frac{uv}{h}=\min(u,v).
\]
So the path has the Wiener law on this cell. Concatenating independently over
the deterministic cells, we obtain jointly Gaussian evaluations at
any finite set of times, with covariance $\min(s,t)$, including across
cells. Since evaluations at rational times determine a Borel law on
continuous paths, this identifies the law of the concatenation. The resulting
pairs are independent across rows, and the $B$ paths are independent
of the $D$ paths. Since \eqref{u:eq:append-cell} preserves every earlier
initial segment, we find by induction that the concatenation is the
reconstructed terminal state. The assertions also hold when unused bridges
are retained.
\end{proof}

\section{Spin laws, histories, and filtrations}\label{u:sec:proposal-laws}

\begin{definition}[The two spin laws and the replay space]\label{u:def:laws}
On the source, we define
\begin{equation}\label{u:eq:spin-link}
 L=\prod_{i=1}^n(1+x_i m_i(V)),\qquad
 \Lambda^{\rm src}=L\mu^{\rm src}_0.
\end{equation}
The unlinked replay law $\mu^{\rm rep}_0$ is the product of
$\mathcal G_n$, independent $C_j\sim N(0,I_n)$, all initial $B^0$
paths and bridges, the $n$ independent $D$ paths, and $\Unif_n$.
Its paths are calculated by \eqref{u:eq:replay-endpoint} and
\eqref{u:eq:append-cell}. We define
$\Lambda^{\rm rep}=L^{\rm rep}\mu^{\rm rep}_0$ using the terminal
fields in replay. We write $\Omega^{\rm rep}$ for this product space,
with the independent antisymmetric variable adjoined when using
primitive disorder. The output law $Q_V$ is the $(g,x)$ marginal of
$\Lambda^{\rm rep}$ on this primitive lift.
\end{definition}

For each fixed choice of all inputs except the spin, we have
$2^{-n}\sum_x L=1$. This implies that both linked laws are probability measures,
and that linking does not change the marginals obtained by forgetting the spin,
including the marginal of the Gaussian matrix. Conditional on the paths,
the terminal spins are independent with masses
$(1+x_i m_i(V))/2$. Equivalently, we may draw the spins using independent
uniform $[0,1]$ thresholds, with the same integrals in either
realization.

Let $\Omega^{\rm row}_V=\{-1,1\}\times\mathcal C_V^3$ be the
canonical Borel space of histories of $(x,Y,B,D)$ for one row through
$V$.
We continue to write $\nu_V$ for the one-row law of
\cref{u:sec:scalar} mapped into this space by the reflected pair and
the Picard solution. From $(\Omega^{\rm row}_V)^n$ and independent
ancillary coordinates, we reconstruct the source as follows. We
restrict $B$ to cell zero and recover, for $j\ge1$,
\begin{equation}\label{u:eq:history-reconstruction}
 \begin{aligned}
 e_{j,i}&=\frac{B_i(t_{j+1})-B_i(t_j)}{\sqrt{h_j}},\\
 \mathrm{br}_{j,i}(u)&=B_i(t_j+u)-B_i(t_j)
          -\frac{u}{h_j}\bigl(B_i(t_{j+1})-B_i(t_j)\bigr).
 \end{aligned}
\end{equation}
The ancillary coordinates are $d_j,v_j,g^0$, the antisymmetric variable
$A$ used in the final lift, and the unused bridges; they are independent, with law
\begin{equation}\label{u:eq:ancillary-law}
 \begin{aligned}
 \mu^{\rm anc}={}&N(0,2/n)^{\otimes N}\otimes N(0,I_n)^{\otimes N}
                  \otimes\gamma_n\otimes\gamma_n^{\rm anti}\\
 &\hspace{12mm}\otimes
       \bigotimes_{j\notin\{1,\ldots,N\}}
                     \mathrm{Br}_{h_j}^{\otimes n}.
 \end{aligned}
\end{equation}
The last factor is omitted if unused bridges are not retained, and
$\gamma_n^{\rm anti}$ is specified below. We write $\mathcal A$ for
the ancillary product space. Since $W_0=W(g^0)$, we retain the whole primitive matrix $g^0$,
including its unused antisymmetric coordinates.

The endpoints and bridges in \eqref{u:eq:history-reconstruction} are
independent, since their evaluations at any finite set of times are jointly
Gaussian with zero cross covariance. We have therefore identified
the product law on the source space. Using the causal
inverse of the reflected pair in \cref{u:sec:scalar}, we identify the
innovation and the scalar field in this reconstruction. All these maps are Borel
under the conventions of \cref{u:sec:setup}. Indeed, evaluations at finitely many times,
restrictions of paths, integration against weights of bounded variation, and
the Picard solution map are Borel,
and the frame rule uses finitely many measurable zero tests and
normalizations.

We record the laws in the following diagram, in which every arrow except the
density arrow denotes an equality of laws under a measurable pushforward. The same
top arrows, applied before linking, identify the unlinked laws. The map $F$
retains the initial data and $D$ and replaces the source noises by
$(W,C)$ from \eqref{u:eq:assembled-matrix}.
\begin{equation}\label{u:eq:law-diagram}
 \begin{gathered}
 \bigl((\Omega^{\rm row}_V)^n\times\mathcal A,
       \nu_V^{\otimes n}\otimes\mu^{\rm anc}\bigr)
 \xrightarrow{\ \operatorname{Rec}\ }
       (\Omega^{\rm src},\Lambda^{\rm src})
 \xrightarrow{\ F\ }
       (\Omega^{\rm rep},\Lambda^{\rm rep})
 \\[4pt]
 \begin{array}{ccc}
 (\Omega^{\rm rep},R_D=E_D\mu^{\rm rep}_0)
 &\xrightarrow{\ d\Lambda^{\rm rep}/dR_D=L^{\rm rep}/E_D\ }&
 (\Omega^{\rm rep},\Lambda^{\rm rep})\\
 \downarrow\scriptstyle(g,x)&&\downarrow\scriptstyle(g,x)\\
 (\R^{n\times n}\times\Sigma_n,\gamma_n\otimes\Unif_n)&&Q_V.
 \end{array}
 \end{gathered}
\end{equation}
Here $\operatorname{Rec}$ is the reconstruction just described,
using \eqref{u:eq:history-reconstruction} and the reflected pair.
Along the density arrow, we change measure on the same replay space; $E_D$
and $R_D$ are constructed in \cref{u:sec:reference}. In particular,
the top row implies the identity for the complete histories
\begin{equation}\label{u:eq:row-history}
 \bigl((x_i,Y_i,B_i,D_i)_{i=1}^n\bigr)_*\Lambda^{\rm rep}
       =\nu_V^{\otimes n}.
\end{equation}

To check the linked identities, let $\nu_V^0$ be the canonical
history law before the one-row link. By the factorization into
endpoints and bridges and the inverse of the reflection, we have
\[
 \operatorname{Rec}_*\bigl((\nu_V^0)^{\otimes n}
                          \otimes\mu^{\rm anc}\bigr)
       =\mu^{\rm src}_0.
\]
On a supporting set of full measure, the recomputed field agrees with
the given $Y$. For every bounded Borel test function $f$ on the source space, we have
\[
 \begin{aligned}
 &\int f\bigl(\operatorname{Rec}(w,a)\bigr)
      \prod_{i=1}^n\bigl(1+x_i m(V,Y_i(V))\bigr)
      \dd\bigl((\nu_V^0)^{\otimes n}\otimes\mu^{\rm anc}\bigr)(w,a)\\
 &\hspace{35mm}=\int fL\dd\mu^{\rm src}_0.
 \end{aligned}
\]
This gives the law on the source space in the top row, with every
ancillary coordinate retained. Repeating this calculation of densities and maps for $F$,
and using the pointwise replay identity in \cref{u:lem:assembly}, we obtain the
pushforward identity for the linked replay law. The matrix and completion vectors are computed
from these rows and ancillary variables; they need not be independent
of the rows.

For the primitive $g$ in \eqref{u:eq:law-diagram}, we retain the
antisymmetric Gaussian variable $A$, independently of all the other
inputs.
Its entries $A_{ij}$, $i<j$, are independent $N(0,1/2)$, with
$A_{ji}=-A_{ij}$ and $A_{ii}=0$. We define
\[
 g=\sqrt{n/2}\,W+A.
\]
Then $g\sim\gamma_n$ and $W(g)=W$, as in \cref{u:sec:primitive}.
Since replay reads only $W$, this gives the proposal with primitive input.
The proposal and the target have the same conditional law of $A$ given
$(W,x)$. Their joint relative entropy on primitive matrices equals
their relative entropy on physical matrices. We use this lift whenever
we condition on $g$.

\paragraph{Filtrations.}
Whether a Brownian assertion may be used depends on the law and on the
information revealed at time zero. The following table records the
distinctions needed in the later proofs. In the last row,
$\xi_i(v)=x_i\chi(v,Y_i(v))/(1+x_i m(v,Y_i(v)))$, as in
\cref{u:sec:reference}.
\begin{center}
\small
\renewcommand{\arraystretch}{1.14}
\begin{tabular}{@{}>{\raggedright\arraybackslash}p{0.29\textwidth}
 >{\raggedright\arraybackslash}p{0.39\textwidth}
 >{\raggedright\arraybackslash}p{0.25\textwidth}@{}}
\toprule
Law and filtration & Brownian or independence assertion & Use \\
\midrule
Scalar Wiener pair, with the spin-free pair past
& $(M,\widetilde M)$ and its deterministic orthogonal transform
  $(B,D)$ are Wiener pairs.
& Scalar martingales, brackets, and path estimates. \\[3pt]
Unlinked row, with an independent uniform sign revealed initially
& $(B,D)$ remains a Wiener pair in the filtration including this sign.
& Signed integrals and fixed-sign It\^o identities. \\[3pt]
Linked row $\nu_V$, with the sign revealed initially
& The entire path marginal is unchanged; $B,D$ need not be Brownian
  in the filtration including the sign.
& Posterior moments, path-only estimates, and the previously defined
  signed integrals. \\[3pt]
Product of complete linked rows and independent ancillary variables
& The complete rows are independent before conditioning on $W$.
& Independent-row concentration estimates. \\[3pt]
Unlinked replay, with $(W,C,B^0,\mathrm{br},x)$ revealed initially
& $D$ is Brownian and independent of the static inputs. No Brownian
  property for $B$ is asserted in this filtration.
& Construction and normalization of $R_D$. \\[3pt]
Linked source natural filtration, including $x$, excluding the final
matrix and unused matrix noises
& The compensated vectors
  $B-\int\rho\xi\dd v$ and $D-\int\sigma\xi\dd v$ are Brownian.
& The mean information identity by bounded-integrand Girsanov. \\
\bottomrule
\end{tabular}
\end{center}
Once the terminal sign or the assembled matrix is revealed, a Brownian
assertion must be verified separately in the resulting filtration.
We transport all scalar expectations at fixed $V$ and all Gram tails below
by \eqref{u:eq:row-history}. We use the bound $L\le2^n$ only for
null sets at fixed $n$ or for integral versions, as in \cref{u:sec:setup}.

\section{Population entries}\label{u:sec:proposal-entries}

For a single linked row, let
$e_j=(B(t_{j+1})-B(t_j))/\sqrt{h_j}$ and
$\mathcal T=(\mathcal T_j)_{j=1}^N$, $e=(e_j)_{j=1}^N$.
By \cref{u:lem:cross}, the identities for the reflected pair in
\cref{u:sec:scalar}, and posterior projection, we have
\begin{equation}\label{u:eq:population-entries}
 \begin{gathered}
 \E\mathcal T\mathcal T^{\mathsf T}=I_N,\qquad
 \E ee^{\mathsf T}=I_N,\qquad
 \E\mathcal T x=s,\qquad \E ex=b,\\
 H_0=\E\mathcal T e^{\mathsf T},\qquad
 (H_0)_{ij}=0\quad(i\le j).
 \end{gathered}
\end{equation}
Indeed, disjoint martingale increments are orthogonal.
The identity $\E\mathcal T_jx=s_j$ follows from the variance identity
in \cref{u:prop:scalar}, and the reflected cross moment gives
$\E e_jx=b_j$. For $i\le j$, $\mathcal T_i$ is measurable before
the fresh $B$ cell defining $e_j$, so the path-only moment is zero.
We retain the entire lower block $H_0$. By Bessel's inequality in the
$L^2$ space for one row, we have
\begin{equation}\label{u:eq:population-bounds}
 \norm{s},\norm{b},\norm{H_0}\le1,\qquad
 0\le e_h=s^{\mathsf T}b\le1.
\end{equation}
The row means vanish by zero-field symmetry. Note that these calculations do not use
independence of the entries within a row.

\begin{lemma}[Adjacent-cell quadrature and stopping]\label{u:lem:quad}
If $\max_{j\le N}h_j\le1/8$, then
\begin{equation}\label{u:eq:adjacent-quadrature}
 |e_h-e(V)|\le C_\beta h.
\end{equation}
For $h=n^{-\ex a}$ and $\ex L>0$, define
\[
 N_{\ex L}+1=\min\{j:t_j\ge\ex L\log n\},\qquad
 V_{\ex L}=t_{N_{\ex L}+1}.
\]
If $\ex a>\ex L/4$, then eventually
\begin{equation}\label{u:eq:growing-stopping}
 \begin{gathered}
 N_{\ex L}+1\le8n^{\ex a},\qquad
 \ex L\log n\le V_{\ex L}
                \le\ex L\log n+n^{-\ex a+\ex L/4},\\
 \max_{j<N_{\ex L}+1}h_j\le n^{-\ex a+\ex L/4}.
 \end{gathered}
\end{equation}
There are $N_{\ex L}+1$ cells and exactly $N_{\ex L}$ matrix queries.
If additionally $\ex a<1$, then eventually $N_{\ex L}+1\ge2$ and
$N_{\ex L}<n$.
\end{lemma}

\begin{proof}
We compare the normalization from the preceding cell with $r$ on the
next cell. Using the relative variation bounds of \cref{u:prop:scalar}, we
obtain an integrable quadrature error. For the stopping estimate, we
sum the decrements of $e^{-t_j/4}$.

For $v\in[t_j,t_{j+1}]$ and $u\in[t_{j-1},t_j]$, the relative
variation bounds, monotonicity, and
$h_j/h_{j-1}=e^{h_{j-1}/4}$ imply that
\begin{equation}\label{u:eq:adjacent-ratio}
 1-h_j/8\le\frac{s_j}{\sqrt{h_jr(v)}}\le1+4h_j.
\end{equation}
For the upper bound, we have $r(v)\ge(1-4h_j)r(u)$ and
$(1-4h_j)^{-1/2}\le1+4h_j$ for $h_j\le1/8$.
For the lower bound, we use $r(u)\ge r(v)$ and
$e^{-h_{j-1}/8}\ge1-h_j/8$. It follows that
\[
 \left|\frac{s_j}{\sqrt{h_j}}-\sqrt{r(v)}\right|
       \le4h_j\sqrt{r(v)}.
\]
With $q_e(v)=\beta r(v)p(v)$, we obtain
\begin{equation}\label{u:eq:quadrature-error}
 |e_h-e(V)|\le h+4h\int_0^V e^{v/4}q_e(v)\dd v.
\end{equation}
The first term $h$ bounds the contribution of the omitted auxiliary cell. The integral is
bounded uniformly in $V$. Indeed, on the finite interval before the endpoint of the support, we have
$q_e\le\beta$, and after it, using $r\le p$ and \cref{u:lem:tails},
we find that the integrand decays exponentially. This proves
\eqref{u:eq:adjacent-quadrature}, with, for example,
\[
 C_\beta=1+4\beta\int_0^\infty e^{v/4}p(v)r(v)\dd v<\infty.
\]

For $h_j\le1$, we have $1-e^{-h_j/4}\ge h_j/8$, and hence
\[
 e^{-t_j/4}-e^{-t_{j+1}/4}\ge h/8.
\]
Summing up to the first reaching time gives
$(N_{\ex L}+1)h\le8$. Every left endpoint that is used lies below
$\ex L\log n$, which proves the bounds on the number of cells and on the overshoot in
\eqref{u:eq:growing-stopping}. The reaching index exists because
$t_j\ge jh$. Since $\ex L\log n>0$, it is at least one.
If additionally $\ex a<1$, it is eventually at least two, because
$t_1=n^{-\ex a}\to0<\ex L\log n$ and
$\ex a>\ex L/4>0$. Further,
\eqref{u:eq:growing-stopping} implies that
\[
 N_{\ex L}<N_{\ex L}+1\le8n^{\ex a}<n
 \qquad\text{once }n^{1-\ex a}>8.
\]
In every fine or coarse use, the fine-energy conditions in
\cref{u:prop:fine-kl} supply $0<\ex k_{\rm f}<1-\ex a$, so
$\ex a<1$; in particular, $\ex a=1/2$ in the exponent table and
the range family of \cref{u:sec:parameters,u:sec:budget}.
We never use $\ex L=0$. All scalar references are
evaluated at $V_{\ex L}$, including its overshoot. The estimate needs neither an extra cell nor a quadrature error
bounded by the maximal duration of a cell.
\end{proof}

We use the same sequence \eqref{u:eq:growing-schedule} for both
$\ex L=\ex b$ and $\ex L=\ex f$, where $0<\ex b<\ex f$,
and the two replays use identical chronological inputs through
$V_{\ex b}$. By induction on the queries, every scalar
initial segment through $V_{\ex b}$ is the same in the two replays. The auxiliary cells agree. If the scalar paths agree through
the preceding cell, then the targets agree, the same remaining frame
and matrix give the same \eqref{u:eq:replay-endpoint}, and the same
bridge appends the same next cell. This argument covers the
deterministic fallback as well. At query $j$, the number of remaining queries
is $N_{\ex L}-j$; this does not change the chronological schedule.
The terminal spins drawn at the fine horizon are never fed back into
the continuation.
